\section{问题二：运输任务与共享电池联合调度}
\label{sec:q2}

\subsection{问题分析与求解思路}
\subsubsection{组批、访问顺序与配送时限的耦合}
本问允许一个架次依次访问多个服务区。改变箱组不仅改变装载与飞行能耗，也改变交接时间；改变访问顺序则影响逐箱交付时刻。因此，物理上能够完成的一组任务，还可能因关键物资送达过晚而无法作为系统方案。本文以第\ref{sec:foundation}章完整任务属性为基础，将组批、机型与时序联合评价，而不以总距离或平均送达时间代替箱级时限。

\subsubsection{飞机释放与电池恢复的不同步}
同一任务开始时必须同时具备一架兼容运输机和一组满电电池。飞机返航后可以换用另一组电池，而原电池继续恢复。故单独优化飞机任务顺序不足以保证日程可执行。有限资源调度需要显式处理占用与接续\cite{pinedo2016scheduling}；本题进一步以两种不同长度的区间描述同一架次，见式\eqref{eq:common_occupancy}。

求解采用两条相互独立的路线：通过任务重构与真实资源重放寻找可行上界，通过任务列松弛和全覆盖下界材料评价质量。先给出可以执行的具体安排，再说明其最优性范围，最后检验扰动下原安排与可修复安排的差异。

\subsection{多架次联合调度模型}
\subsubsection{运输任务选择与货箱完整配送}
设物理可行任务集为$\mathcal C$。任务$r$包含机型$m_r$、箱组$\mathcal B_r$、路线$\pi_r$、时长$p_r$、能耗$E_r$、恢复时长$c_r$及交付偏移$\delta_{rb}$。变量$x_r\in\{0,1\}$表示选用该任务，则
\begin{equation}
 \sum_{r:b\in\mathcal B_r}x_r=1,\qquad b\in\mathcal B.
\label{eq:q2_cover}
\end{equation}
候选生成已检查质量、体积、完整飞行能量与路径一致性。实现采用按需生成与邻域重构，不要求在排程前显式存储全部可能任务。

\subsubsection{飞机、电池指派与时间冲突约束}
令$\mathcal U_m,\mathcal K_m$为机型$m$的实体飞机和电池集合，二元变量$a_{ru},b_{rk}$表示任务指派，开始时刻$s_r\ge0$。约束为
\begin{equation}
 \sum_{u\in\mathcal U_{m_r}}a_{ru}=x_r,\qquad
 \sum_{k\in\mathcal K_{m_r}}b_{rk}=x_r.
\label{eq:q2_assignment}
\end{equation}
同一实体不能承担重叠任务：
\begin{equation}
\begin{aligned}
 a_{ru}=a_{r'u}=1&\Rightarrow I_r^U\cap I_{r'}^U=\varnothing,\\
 b_{rk}=b_{r'k}=1&\Rightarrow I_r^B\cap I_{r'}^B=\varnothing.
\end{aligned}
\label{eq:q2_disjoint}
\end{equation}
飞机区间结束于返航，电池区间结束于充满。上述析取可用任务先后变量线性化，也可由后述事件排程直接构造满足约束的解。本题未给出充电器数量约束，故不另设共享充电队列，但不能提前复用正在恢复的同一电池。

\subsubsection{硬截止、期望送达与目标优先关系}
被选任务携带货箱$b$时，$C_b=s_r+\delta_{rb}$。医疗物资及指定首批保障货箱须满足相应硬截止$C_b\le h_b$，不能用节能收益补偿违约。期望送达$d_b$与优先权重$\omega_b$用于
\begin{equation}
 L_b=\max(0,C_b-d_b),\qquad L_w=\sum_b\omega_bL_b.
\label{eq:q2_tardiness}
\end{equation}
在硬约束可行的基础上采用
\begin{equation}
 \min_{\mathrm{lex}}(L_w,T_{\max},E,N),\quad
 T_{\max}=\max_{r:x_r=1}(s_r+p_r),\quad E=\sum_rx_rE_r,\quad N=\sum_rx_r.
\label{eq:q2_objective}
\end{equation}
先保障及时交付，再缩短最后返航，随后比较能耗和架次。正式方案$L_w=0$，结合非负性可证明第一层目标已达最优，但这不能推出第二层工期已全局最优。

\subsubsection{任务启用与资源先后的显式连接}
式\eqref{eq:q2_disjoint}表达同一资源上的析取关系。为展示它如何转化为优化约束，对不同任务$r,r'$和同型飞机$u$，引入$z^U_{rr'u}\in\{0,1\}$，表示两项任务同时分给$u$时是否由$r$先执行。在有限候选模型中，取覆盖所考虑时间域及最长占用长度的常数$M_U$，可写为
\begin{equation}
\begin{aligned}
s_{r'}&\ge s_r+p_r-M_U(1-z^U_{rr'u})-M_U(2-a_{ru}-a_{r'u}),\\
s_r&\ge s_{r'}+p_{r'}-M_Uz^U_{rr'u}-M_U(2-a_{ru}-a_{r'u}).
\end{aligned}
\label{eq:q2_plane_disjunction}
\end{equation}
当两项任务都使用$u$时，两式只保留选定方向的先后约束；至少一项没有分给$u$时，两个不等式均被解除。候选时间域可由已知可行日程给定，$M_U$应由实际上界推得，而不是使用没有量纲依据的极大数。电池关系采用独立先后变量，并将飞机占用长度$p_r$替换为$p_r+c_r$，同时为充电结束时刻保留相应上界。

这一显式表达说明：飞机与电池可以分别有不同的前后任务，同一开始时刻$s_r$则必须满足两条资源链的释放条件。它与事件排程对应同一个物理模型。前者方便解释完整约束，后者在大批候选之间构造可行日程更直接；使用事件排程并不意味着忽略了两类资源冲突。

\subsubsection{分层目标与搜索评分的分工}
硬截止先作为可行性约束，随后按$(L_w,T_{\max},E,N)$比较候选。该顺序可由逐层固定最优值实现，也可直接用于可行方案档案的字典序比较。其含义是，前一层取得改善时，不用后一层的变化抵消它；所有指标仍分别保留原始单位和数值。

搜索阶段采用的标量评分用于引导探索，允许暂时接受较差当前状态以离开局部结构。正式最好解从经过完整检查的可行档案中选择，始终采用既定目标顺序。因此，搜索评分中的能耗系数并不是将1 kWh与若干秒等价为社会成本。论文通过具体方案和独立指标呈现取舍，使算法参数与救援评价保持分工。

少飞一个架次也未必缩短工期。拆分任务可能降低单次载荷、提高并行性，却增加准备、交接和资源接续；合并任务则可能节省固定成本，却把关键箱子的交付推迟。候选变化必须经过完整任务计算和双资源时间表生成，才能评价系统效应。


\subsection{调度方案的构造与改进}
\subsubsection{从任务顺序生成执行时间表}
一个候选编码为任务序列$(r_1,\ldots,r_N)$，每项保存机型、原始箱号和访问顺序。对每种机型分别维护飞机、电池最早可用时刻。当前任务选择的资源分别在$t_u,t_k$可用，则
\begin{equation}
 s_r=\max(t_u,t_k),\qquad
 t_u\leftarrow s_r+p_r,\qquad t_k\leftarrow s_r+p_r+c_r.
\label{eq:q2_decoder}
\end{equation}
这一步称为事件排程或解码：它把任务顺序转换为完整资源时间表。对给定序列准确计算时刻，不等于求出了原组合问题的全局最优日程。修改箱组、机型或路线后，必须重算物理属性、充电时间及箱级交付，再执行解码。

\subsubsection{货箱移除、重新插入与局部调整}
本文采用自适应大邻域搜索的破坏--修复框架\cite{ropke2006alns}，把标准框架中的操作对象具体化为货箱、访问顺序和飞机--电池任务链。一次破坏使用随机移除3--7箱、移除一至两项完整任务、移除空间相关货箱，或从当前平均负载最重机型的长任务中移除货箱。删除服务区后也重新计算航段；山区几何不保证直接删点后的航段代价更小，不能跳过物理核验。

修复时，对每个待插入货箱尝试现有任务及单独新建任务，允许重新选择A/B/C机型。单个新增服务区遍历原访问顺序中的插入位置；需要重构多个服务区时，至多5个不同服务区枚举排列，更多服务区仅保留最近邻顺序及反序。后一规则是搜索筛选，不是题目对访问数量的限制。先按平均机型工作量、能耗和时限必要条件排序，保留前10项及每种机型最好的1项，再对去重后的候选遍历该机型任务链中的插入位置，执行式\eqref{eq:q2_decoder}的完整排程。因而这里没有声称穷举全部插入组合。

设货箱$b$的可行插入评分从小到大为$F_{b1},F_{b2},\ldots$。贪心修复选择当前最低评分；$k$阶后悔修复优先处理
\begin{equation}
 R_k(b)=\sum_{j=2}^{k}(F_{bj}-F_{b1}),\qquad k\in\{2,3\}
 \label{eq:q2_regret}
\end{equation}
较大的货箱，即优先安排缺少好替代位置的货箱。可用位置不足$k$个时重复最后一项评分。每轮修复后再进行20次局部提议：同机型任务换序、任务改型、跨任务换箱和单箱迁移的选择概率依次为0.40、0.18、0.22、0.20；全部提议仍需重算物理属性和真实资源日程。

\textbf{物理等价货箱的原始编号匹配。}同区同类箱子可能质量、体积相同而截止要求不同，求解器不能只关心总质量。实现先依据当前日程生成逐箱交付事件，再在同一服务区、同一物资类型的组内，把原始箱号按“硬截止、期望送达、箱号”排序，把交付事件按实际时刻排序，依次匹配。无硬截止记为排序中的正无穷。当前输入中各匹配组的质量和体积一致，故交换身份不改变任务箱数、载荷、装卸时间和能耗；每箱的原始时限及权重始终随原始箱号保留。匹配后重新排程并逐箱验收，而不直接宣称该排序解决了任意加权迟到问题。若输入不满足组内物理等价，则必须细分组，不能沿用此交换。

\subsubsection{搜索接受规则与独立方案核验}
最终方案始终按式\eqref{eq:q2_objective}比较；为在搜索过程中探索不同邻域，程序另采用标量代理评分。记$n_H$为硬迟到箱数、$L_H$为硬迟到秒数之和，则工期搜索模式下
\begin{equation}
\begin{split}
 F(S)={}&2\times10^6n_H+2000L_H+100L_w+T_{\max}\\
        &+1.5E+2N+2\times10^{-4}\sum_bC_b .
\end{split}
\label{eq:q2_search_score}
\end{equation}
式中时间用秒、能耗用kWh；系数是代码固定的搜索导航设置，不是将不同量纲变为等价救援成本。硬迟到候选即使参与探索，也不能进入正式可行方案档案。

令迭代上限为$K$、当前轮为$k$。大邻域和局部提议分别使用温度
\begin{equation}
 \Theta_k=55(1-k/K)^2+1,\qquad
 \theta_k=25(1-k/K)^2+0.2.
 \label{eq:q2_temperature}
\end{equation}
评分更低的候选直接接受，否则以$\exp[-(F(S')-F(S))/\Theta_k]$的概率接受；局部操作将$\Theta_k$替换为$\theta_k$。这使当前搜索状态能够暂时变差，但最好可行方案仅在词典序严格改善时更新。

破坏与修复算子的初始权重为1，按权重比例抽取。某大邻域操作得到新的最好可行解、只改善当前代理评分、或仅被接受时，奖励依次为20、6、1；未接受为0。每25轮对使用过的算子$o$更新
\begin{equation}
 w_o\leftarrow\max\{0.15,\;0.8w_o+0.2\bar r_o\},
 \label{eq:q2_operator_weight}
\end{equation}
其中$\bar r_o$是本周期每次使用的平均奖励；未使用时保持其权重。每80轮把当前状态重置为最好可行解，属于搜索回退而非从零重启。达到迭代上限或墙钟上限之一即停止。

\begin{table}[htbp]\centering\small
\caption{任务重构与双资源搜索的一次完整迭代}\label{tab:q2_algorithm}
\begin{tabularx}{\linewidth}{cX}\toprule
步骤 & 输入、操作及更新对象\\\midrule
1 & 读取完整任务序列，重算物理属性、原始箱号交付及飞机--电池日程；保存当前状态与最好可行方案。\\
2 & 依据权重选择破坏、修复算子，移除相应货箱；按候选筛选规则生成插入位置、机型与访问顺序。\\
3 & 用贪心或后悔规则修复，匹配物理等价原始箱号，并完整核验容量、能量、时限和资源接续。\\
4 & 按代理评分和降温规则更新当前状态；仅对可行且词典序更好的解更新最优档案并记改进日志。\\
5 & 提出20次局部调整，再次核验与接受；按25轮和80轮周期分别更新算子权重和回退当前状态。\\
6 & 达到轮数或时间限制后，输出最好可行方案；独立程序从原始参数重新核验，而不复用主评价器的中间结果。\\\bottomrule
\end{tabularx}
\end{table}

\textbf{正式方案与新增试验的来历。}已保存的输入基线含23架次，工期5693.231489\,s、能耗66.225670\,kWh；选定记录标记为种子9201的后续候选。核对前后决策，可识别S001两项C型任务的一只食品箱与一只卫生用品箱交换，净节能0.011211\,kWh，工期不变。当前归档可以逐项重放这两份方案，但未保存这次历史搜索的完整逐轮日志和实际停止轮数，不能据此编造从空任务集开始的收敛曲线。后文20次试验采用另一共同受扰起点和150轮预算，1\%、2\%能耗备选各250轮；这些实际设置与历史最好方案明确分开。标准方法、候选来源和复现范围见附录\ref{app:source_provenance}。

独立验证逐项重建节点、货箱、航段、能耗、充电和交付，并检查每个飞机和电池区间。正文表\ref{tab:q2_schedule}、箱号附表与甘特图均由同一份23架次、80箱交付记录生成，原始决策列表顺序不再与按执行时刻排列的编号混用。


\subsubsection{插入机会与资源瓶颈的配合}
后悔插入比较一个待配送货箱的最好插入位置与若干替代位置。可用替代方案少、最好位置又明显便宜的货箱，应优先安排，避免其可行插入机会被其他箱子占用。贪心修复强调当前最小增量，后悔修复强调未来替代成本。两种策略从不同角度组织同一个候选池，不需要人为保证某一策略始终最好。

Pisinger与Ropke用统一框架处理多种车辆路径约束，展示了多类移除和插入策略的组合方式\cite{pisinger2007}。本题将插入代价进一步落实到机型重选和双资源接续：同一箱子的目的地相同，插入不同任务后，受影响的载荷序列、交接偏移与充电长度仍可能不同。几何缓存减少重复地形查询，候选筛选减少完整解码次数，而最终入档保持真实物理与时限核验。

自适应权重分配的是计算机会。Turkeš等对自适应大邻域搜索的元分析表明，适应层的增益与算子选择、实现和比较设置有关\cite{turkes2021}。因此本文列明权重更新与预算，用规定起点下的重复运行描述本实例表现，再用独立下界评价保存方案质量。关键任务移除的针对性来自负载和接续关系，而不来自未经验证的“智能优化”标签。

\subsubsection{一次箱组调整怎样传播到完整日程}
考虑将货箱$b$从任务$r$移入$r'$。首先更新两项箱组，删除不再需要访问的服务区，并尝试目标目的地的允许插入位置。对保留路线和兼容机型，重新检查质量、体积与返航能量；每站交接偏移和返航SOC也随箱数及载荷变化重新计算。超限候选在完整排程前即被排除。

随后将更新后的任务放入相应机型序列，按最早可用飞机和满电电池递推开始时刻。原任务充电时间变化会传到后续用同一电池的任务，继而可能再影响后继飞机链。最后把实际交付事件与原始箱号的期限属性相匹配，检查80箱的唯一配送、硬截止和加权迟到。局部移动一个箱子，由此形成“箱组—物理成本—资源释放—后继开始—交付时限”的完整反馈。

上述组织把计算代价分为物理评价、插入筛选、资源排程与箱级检查四部分。无关机型的已验证属性可以保留，受影响机型的后续事件则必须重放。这既利用局部修改的结构，也避免仅比较变化任务自身的时间而遗漏全局影响。


\subsection{调度结果与质量评价}
\subsubsection{运输路线及飞机—电池执行安排}
正式方案共23架次，A/B/C为11/6/6，使用8架运输机和14组电池，最后返航5693.231489\,s，总能耗66.214460\,kWh。表\ref{tab:q2_schedule}按实际架次编号列出资源与时刻；完整箱号在附录\ref{app:full_schedules}。图\ref{fig:q2_gantt}中同一任务在飞机层和电池层保持对应。
\begin{longtable}{@{}lclllrr@{}}
\caption{问题二正式23架次的资源指派与执行时刻}\label{tab:q2_schedule}\\
\toprule
架次 & 机型 & 访问顺序 & 飞机 & 电池 & 开始/min & 返航/min\\
\midrule
\endfirsthead
\toprule
架次 & 机型 & 访问顺序 & 飞机 & 电池 & 开始/min & 返航/min\\
\midrule
\endhead
Q2-R-01 & A & S001 & U01 & A-BAT01 & 0.000 & 20.848\\
Q2-R-02 & A & S004 & U02 & A-BAT02 & 0.000 & 36.313\\
Q2-R-03 & A & S013 & U03 & A-BAT03 & 0.000 & 26.869\\
Q2-R-04 & A & S011$\to$S004 & U04 & A-BAT04 & 0.000 & 41.068\\
Q2-R-05 & B & S010 & U05 & B-BAT01 & 0.000 & 25.923\\
Q2-R-06 & B & S014 & U06 & B-BAT02 & 0.000 & 31.163\\
Q2-R-07 & C & S001 & U07 & C-BAT01 & 0.000 & 27.103\\
Q2-R-08 & C & S001 & U08 & C-BAT02 & 0.000 & 23.803\\
Q2-R-09 & A & S009$\to$S002$\to$S013 & U01 & A-BAT05 & 20.848 & 65.440\\
Q2-R-10 & C & S006 & U08 & C-BAT03 & 23.803 & 49.829\\
Q2-R-11 & B & S012 & U05 & B-BAT03 & 25.923 & 57.969\\
Q2-R-12 & A & S007 & U03 & A-BAT06 & 26.869 & 55.659\\
Q2-R-13 & C & S002 & U07 & C-BAT04 & 27.103 & 63.333\\
Q2-R-14 & B & S004 & U06 & B-BAT04 & 31.163 & 63.190\\
Q2-R-15 & A & S007 & U02 & A-BAT01 & 36.313 & 65.104\\
Q2-R-16 & A & S003$\to$S015 & U04 & A-BAT03 & 44.954 & 88.343\\
Q2-R-17 & C & S007$\to$S003 & U08 & C-BAT02 & 49.829 & 94.722\\
Q2-R-18 & B & S008 & U05 & B-BAT01 & 57.969 & 90.666\\
Q2-R-19 & A & S015 & U03 & A-BAT02 & 60.011 & 92.821\\
Q2-R-20 & B & S008 & U06 & B-BAT02 & 63.190 & 94.887\\
Q2-R-21 & C & S005 & U07 & C-BAT01 & 63.333 & 94.620\\
Q2-R-22 & A & S009 & U02 & A-BAT04 & 65.104 & 93.247\\
Q2-R-23 & A & S011 & U01 & A-BAT06 & 73.166 & 93.533\\
\bottomrule
\end{longtable}

\begin{figure}[htbp]\centering
\includegraphics[width=.98\linewidth]{q2_aircraft.pdf}\par\smallskip
\includegraphics[width=.98\linewidth]{q2_batteries.pdf}
\caption{正式23架次方案的双资源时间表。上图为运输机，下图为电池；浅填充与斜线为充电恢复，虚线为最后返航。所有条块与表\ref{tab:q2_schedule}来自同一份日程。}\label{fig:q2_gantt}
\end{figure}

\subsubsection{箱级时限与返航电量核验}
80箱唯一交付，加权迟到与硬迟到均为0。最后一箱在5283.168387\,s交付，全部运输机在5693.231489\,s返回，相差410.063102\,s，故不能以最后交付替代最后返航。最小硬截止余量202.096058\,s，最低返航SOC为20.097383\%。后者只略高于20\%下限，应理解为名义可行，而非能量误差宽裕。

现有方案曾通过S001两项C型任务的食品箱与卫生用品箱交换，在工期与迟到不变下节省0.011211\,kWh。该调整利用不同载荷组合的能耗差，不改变所有任务的时限可行性。相关箱号以当前正式重放为准。

\subsubsection{机型分工与执行负载}
表\ref{tab:rev_q2_types}按正式日程重新统计运输质量、能耗和工作量。A型11架次承担231 kg，B型6架次承担140 kg，C型6架次承担387 kg。A型以较轻任务提供并行交付能力，C型以较少架次承担较大质量。机型分工由具体任务成本与资源数量共同决定，而非将需求平均拆给八架飞机。
\begin{table}[htbp]\centering\small
\caption{正式运输日程的分机型任务与负载}\label{tab:rev_q2_types}
\begin{tabular}{lrrrrrr}\toprule
机型 & 架次 & 实体数 & 质量/kg & 能耗/kWh & 累计作业/s & 平均负载/s\\\midrule
A & 11 & 4 & 231 & 26.367859 & 21118.826 & 5279.706\\
B & 6 & 2 & 140 & 15.123025 & 11133.183 & 5566.591\\
C & 6 & 2 & 387 & 24.723576 & 11360.487 & 5680.244\\
\bottomrule\end{tabular}
\end{table}


把各类累计作业时间除以实体数量，可得到忽略不可拆任务形态的平均负载下界。C型的该值为5680.243675 s，距正式最后返航仅12.987814 s；B型的平均界为5566.591421 s，其不可拆分配产生更强的离散瓶颈，见下一节。两项统计共同说明，继续改善需要改变接近满负载机型的任务结构，单纯把某一非关键任务提前并不一定改变系统工期。

在共同窗口$[0,T_{\max}]$内，定义飞机$u$的任务占用率及全机群占用率为
\begin{equation}
 U_u=\frac{\sum_{r:a_{ru}=1}p_r}{T_{\max}},\qquad
 U_{\mathrm{fleet}}=\frac{\sum_rp_r}{|\mathcal U|T_{\max}}.
 \label{eq:q2_occupation}
\end{equation}
这里的作业时间包含准备、飞行和交接，因此度量的是任务对实体资源的占用，而非纯空中飞行率。八架飞机的共同窗口占用率为95.76\%。表\ref{tab:rev_q2_util}显示，U06在该窗口上连续作业，两架C型机的占用率均超过99.7\%。这为甘特图中的紧凑安排提供了可复算的量化说明。
\begin{table}[htbp]\centering\small
\caption{共同执行时间窗口内的运输机作业占用}\label{tab:rev_q2_util}
\begin{tabular}{llrrrr}\toprule
实体 & 机型 & 架次 & 作业/min & 非占用/min & 占用率/\%\\\midrule
U01 & A & 3 & 85.808 & 9.080 & 90.43\\
U02 & A & 3 & 93.247 & 1.640 & 98.27\\
U03 & A & 3 & 88.469 & 6.418 & 93.24\\
U04 & A & 2 & 84.457 & 10.430 & 89.01\\
U05 & B & 3 & 90.666 & 4.221 & 95.55\\
U06 & B & 3 & 94.887 & 0.000 & 100.00\\
U07 & C & 3 & 94.620 & 0.268 & 99.72\\
U08 & C & 3 & 94.722 & 0.165 & 99.83\\
\bottomrule\end{tabular}
\end{table}


\subsubsection{交付进度与任务完成层次}
表\ref{tab:rev_q2_milestones}按交接完成事件统计配送进度。30 min完成28箱、268 kg，60 min完成53箱、498 kg，90 min前80箱全部交付。某一早期窗口已交付数量较少，并不表示机群没有工作：首批任务仍可能处于准备、转场或交接阶段，只有完成交接后才能进入交付累计量。
\begin{table}[htbp]\centering\small
\caption{交接完成口径下的累计配送进度}\label{tab:rev_q2_milestones}
\begin{tabular}{lrrrr}\toprule
时刻/min & 已交付箱数 & 质量/kg & 体积/m$^3$ & 箱数进度/\%\\\midrule
15 & 1 & 3 & 0.012 & 1.25\\
30 & 28 & 268 & 0.689 & 35.00\\
45 & 35 & 330 & 0.857 & 43.75\\
60 & 53 & 498 & 1.303 & 66.25\\
75 & 56 & 521 & 1.377 & 70.00\\
90 & 80 & 758 & 2.011 & 100.00\\
\bottomrule\end{tabular}
\end{table}


箱数进度和质量进度反映不同内容。医疗、食品、水和卫生用品的单箱质量不同，故某段时间交付箱数多不必然代表交付质量大。表\ref{tab:rev_delivery_kind}进一步给出各物资类别的交付范围及对应硬时限余量。实际遵守的是每只箱子的原始期限，而不是额外强加“所有医疗箱一定先于所有其他物资”的统一规则。
\begin{table}[htbp]\centering\small
\caption{各物资类别的实际交付时段与硬时限余量}\label{tab:rev_delivery_kind}
\begin{tabular}{lrrrrr}\toprule
物资类型 & 箱数 & 最早交付/min & 最后交付/min & 硬时限箱数 & 最小硬余量/min\\\midrule
医疗物资 & 16 & 14.887 & 86.220 & 16 & 3.368\\
应急食品 & 19 & 15.143 & 88.053 & 0 & --\\
饮用水 & 36 & 15.143 & 88.053 & 15 & 12.889\\
生活卫生用品 & 9 & 21.904 & 86.220 & 0 & --\\
\bottomrule\end{tabular}
\end{table}


从最后一箱交付到全部飞机返航仍有410.063102 s。这段时间属于装备回收，继续占用飞机并产生返航能耗。交付进度和最后返航分别说明物资到位与系统回收两个完成层次，使应急效果和资源周转可以同时核查。

\subsubsection{电池共享的实际接续关系}
23项任务使用14组电池，形成9条再次使用同一电池的接续关系，见表\ref{tab:rev_q2_tight_chains}。其中3条在计算精度内几乎没有充满后的等待。电池池通过不同飞机交替使用已恢复能源支撑多架次，而非一次性分配后永久绑定实体机。
\begin{table}[htbp]\centering\small
\caption{共享电池再次投入任务的接续关系}\label{tab:rev_q2_tight_chains}
\begin{tabular}{lllrrr}\toprule
电池 & 前任务 & 后任务 & 恢复/s & 后项开始/s & 接续余量/s\\\midrule
A-BAT01 & Q2-R-01 & Q2-R-15 & 2121.355 & 2178.800 & 57.445\\
A-BAT02 & Q2-R-02 & Q2-R-19 & 3600.679 & 3600.679 & 0.000\\
A-BAT03 & Q2-R-03 & Q2-R-16 & 2697.242 & 2697.242 & 0.000\\
A-BAT04 & Q2-R-04 & Q2-R-22 & 3895.818 & 3906.224 & 10.406\\
A-BAT06 & Q2-R-12 & Q2-R-23 & 4389.970 & 4389.970 & 0.000\\
B-BAT01 & Q2-R-05 & Q2-R-18 & 3012.098 & 3478.120 & 466.022\\
B-BAT02 & Q2-R-06 & Q2-R-20 & 3463.925 & 3791.400 & 327.475\\
C-BAT01 & Q2-R-07 & Q2-R-21 & 3181.215 & 3799.981 & 618.767\\
C-BAT02 & Q2-R-08 & Q2-R-17 & 2956.347 & 2989.723 & 33.376\\
\bottomrule\end{tabular}
\end{table}


对同电池的相邻任务$r,r'$，定义接续余量
\begin{equation}
 \sigma_{rr'}^B=s_{r'}-(s_r+p_r+c_r).
 \label{eq:q2_battery_slack}
\end{equation}
当恢复时间增加超过该余量，原开始时刻必须调整。接续余量较小的任务是充电压力检验的重点，但局部推迟是否增加最后返航，还取决于另一条资源链何时完成。因此后续实验沿真实依赖传播，不把所有局部延长量直接累加到工期。


\subsubsection{关键资源负载及不同架次数方案比较}
六项B型任务的累计作业时间为11133.182841\,s，分配给两架飞机的平均负载下界为5566.591421\,s。枚举这六个不可拆任务的两机分配，得到更强的离散下界5693.231489\,s，当前日程恰好达到。因而在六项B型箱组、机型、路线固定时，仅换序无法缩短工期；进一步改善须改变任务结构或其他约束，而不是继续调整同一顺序。

原22架次可行见证的工期为7767.058323\,s、能耗66.323024\,kWh；正式23架次方案在两项指标均更小。但该见证不是22架次最优解。本轮短预算稳定性试验还获得6295.573881\,s、71.484859\,kWh的22架次候选，说明22架次方案间也存在较大差异。三者共同支持的是“少一架次不保证更短工期”，而不是“所有22架次方案都必须达到原见证的工期”。

\subsubsection{完成时间下界与可行方案差距}
\label{sec:q2_bound_derivation}
高质量可行解给出上界，但“搜索未改善”不能证明接近全局最优。本文另外采用任务列松弛与互补分支覆盖\cite{barnhart1998branchprice}。下界的关键在于保证每份原问题可行日程都能映射到所求松弛，而不要求松弛解能够反向执行。

\textbf{从实际任务到乐观任务列。}把同一服务区、物资类型、质量及体积的货箱压缩为53个物理类，需求为$d_j$；期限仍留在上界的逐箱验证中，下界不直接施加这些期限。真实航段的时间和能量分别向下取整到毫秒、Wh，再分别作最短路闭包，能量预算向上包络。因此两种闭包不必对应同一条实际航迹，却不会高估实际代价。

在这个删除时限的闭包模型中，可把某服务区全部物资提前到其首次正交付时一起交付，再删除后续重复访问。总装载量、体积和交付数量不变；提前卸货使后续载荷不增，载荷单调能耗与闭包三角不等式保证时间和能耗不增。这个投影论证\textbf{只适用于下界模型}，不能倒过来证明真实山区路线不存在有益的回访。实际定价只要求A、B型不重复访问，C型仍允许回访，因而进一步放松而不排除上述投影。

记乐观列$r$的机型为$g(r)$、交付计数为$a_{jr}$、时间下界为$\underline p_r$（本节统一换算成秒），$n_g$为机型$g$的实体数量。连续主问题为
\begin{equation}
\begin{aligned}
 \min\quad &T\\
 \mathrm{s.t.}\quad&\sum_ra_{jr}x_r=d_j,\quad j=1,\ldots,53,\\
 &\sum_{r:g(r)=g}\underline p_rx_r\le n_gT,\quad g\in\{A,B,C\},\\
 &Bx\le b,\qquad x_r\ge0,\quad T\ge0.
\end{aligned}
\label{eq:q2_relaxed_master}
\end{equation}
根问题的$B$行可为空，节点中则包含有效的整数割和分支条件。任一真实日程的单机累计工作量不超过其最后返航时间，故机型总工作量不超过$n_gT$。这里放松了单架飞机的非抢占接续、共享电池恢复及逐箱时限，故只给下界，不能把分数列解释为配送任务。

\textbf{加强约束与适用条件。}对于货箱类子集$S$，整数覆盖满足
\begin{equation}
 \sum_r\left\lfloor\frac{\sum_{j\in S}a_{jr}}2\right\rfloor x_r
 \le\left\lfloor\frac{\sum_{j\in S}d_j}2\right\rfloor .
 \label{eq:q2_subset_cut}
\end{equation}
其有效性来自整数任务计数和向下取整的次可加性。另在条件$T\le H$下，$H=5465.503$\,s，定义$f_k(p)=\lceil kp/H\rceil-1$。每架飞机上一组任务满足$\sum f_k(p)<(k/H)\sum p\le k$，左侧为整数，所以对$k=2,\ldots,8$有
\begin{equation}
 \sum_{r:g(r)=g}f_k(\underline p_r)x_r\le n_g(k-1).
 \label{eq:q2_packing_cut}
\end{equation}
这些割可排除部分平均工作量合格、却不能装进有限单机时间轴的分数解。它们有$T\le H$的前提，必须把$T>H$的补集一同纳入全局证明。

\textbf{完整定价而不是有限列池的最优值。}对覆盖等式取自由乘子$\pi$，机型工作量和$Bx\le b$分别取非负乘子$\mu,\lambda$。当$\sum_gn_g\mu_g\le1$，且对\emph{所有允许列}有
\begin{equation}
 \bar c_r=\mu_{g(r)}\underline p_r+\lambda^{\mathsf T}B_{\cdot r}
            -\sum_j\pi_ja_{jr}\ge0,
 \label{eq:q2_pricing_rc}
\end{equation}
由弱对偶可得
\begin{equation}
 T\ge\beta=d^{\mathsf T}\pi-b^{\mathsf T}\lambda.
 \label{eq:q2_dual_bound}
\end{equation}
因此，受限列池求得对偶后，还须按机型完整求解$\min_r\bar c_r$。发现负约化成本则加入对应列；没有完成定价时，不把受限主问题目标作为新下界，而保留该区域已有的有效证据。

定价标签保存末节点、质量、体积、能量、时间及相关割状态；A、B型还保存访问集合。两服务区共同配送的分支要求标签保留相应成员状态，只有状态兼容时才能作支配删除，否则可能丢掉负约化成本列。节点中机型$g$对类$j$的总配送量边界为$l_{gj},u_{gj}$时，由覆盖关系再收紧单次停靠的必要上限为$\min\{u_{gj},d_j-\sum_{h\ne g}l_{hj}\}$；这一筛选来自整数可行域，不来自当前列池中是否出现过该组合。

\textbf{互补分支与全域下界。}机型架次数、物理类配送数量、规范化访问次数和两服务区共同配送次数均为整数聚合量。对分数值$v$分成$v\le\lfloor v\rfloor$与$v\ge\lfloor v\rfloor+1$，两子域覆盖该节点的所有规范化真实方案。设完整覆盖的叶域下界为$\beta_\ell$，则包含上述时域补集的全局下界为
\begin{equation}
 LB=\min\left\{H,\min_{\ell\in\mathcal L}\beta_\ell\right\}.
 \label{eq:q2_cover_bound}
\end{equation}
不能取不同叶域下界的最大值，也不能遗漏未解决的子域。不同证明树仅在相同覆盖区域上选择较强的完整子证明；最终仍对整个覆盖取最小值。

正式证据以缩放整数保存乘子，向安全方向修复舍入造成的负约化成本，再对修复后的乘子重新完整定价；分子、分母采用有理数比较。最终下界向下取到微秒，可行上界向上包络。保存的809个分支、810份叶材料和2430次完整机型定价支持
\begin{equation}
 5433.956428\le T^*\le5693.231490\ \mathrm{(s)},\qquad
 \frac{UB-LB}{UB}=4.5541\%.
 \label{eq:q2_bound}
\end{equation}
该确定性区间适用于原始名义实例，不是统计置信区间，也不迁移到改变充电或能耗的新实例。详细叶材料来自已锁定的历史完整认证；本次论文修订核对其数学和数据对应关系，并未重新运行全部2430次定价。附录保留复算入口，正文不以材料数量代替上述有效性链条。


\subsubsection{完整定价怎样把计算结果转为下界}
任务列主问题首先看到的是机型总工作量。两架飞机的工作量除以2给出平均界，但实际任务不能被分成任意小片段，在两条时间轴上的装填可能存在离散不平衡。因此装填割增加了任务长度分布的信息，覆盖约束保证物资数量，分支条件进一步约束整数聚合量，三者共同加强松弛。

另一方面，受限列池可能遗漏便宜任务，直接使用其最优值并不足以形成原问题下界。完整定价检查所有允许列的对偶不等式；Feillet等关于资源约束基本最短路的精确标签算法为此类状态扩展与支配设计提供了方法基础\cite{feillet2004}。本题的定价除了容量和能源，还要保留相关分支状态，使支配删除不会遗漏影响对偶可行性的列。

若全部允许列满足式\eqref{eq:q2_pricing_rc}，则对任一映射到该叶域的实际方案，有
\begin{equation}
\begin{aligned}
 \pi^{\mathsf T}d
 &\le\sum_g\mu_g\sum_{r:g(r)=g}\underline p_rx_r+\lambda^{\mathsf T}Bx\\
 &\le T\sum_gn_g\mu_g+\lambda^{\mathsf T}b
 \le T+\lambda^{\mathsf T}b.
\end{aligned}
\label{eq:q2_weak_dual_chain}
\end{equation}
移项即为$T\ge\pi^{\mathsf T}d-\lambda^{\mathsf T}b$。这个推导说明，完整定价保证的是证书中的对偶可行性，而不只是“列池不再变化”。各叶域分别得到有效界，再通过互补覆盖统一到原问题，才形成最终认证区间。

上界与下界可以沿不同路线持续改善：重构任务和时序降低可行上界，加入有效约束与加强定价提高下界。当前两者相距259.275062 s，量化了名义实例中尚可能存在的工期改进空间。可执行方案与独立界定相结合，使解质量不再仅由一次启发式输出的数值判断。


\subsection{稳健性与敏感性分析}
\label{sec:q2_robust}
\subsubsection{不同随机种子下的结果稳定性}
为避免每次从最终优良方案出发、只重复返回原解而制造“稳定”，先对正式方案进行100次局部扰动提议，保留可行变化，得到一份工期9146.631206\,s的共同可行起点。20个独立随机种子均从该起点出发，各运行150轮搜索，设置120\,s安全上限；全部由轮数限制结束，实测19.919—26.831\,s。每份输出均通过独立物理和完整资源验证。

20次均保持零迟到。工期最小6229.574秒，最大6862.169秒，中位数6457.714秒，样本标准差169.209407\,s，变异系数约2.61\%。架次为22—26，能耗为69.034668—78.336575\,kWh。图\ref{fig:q2_seed}展示全部样本，不只选最好一次。这一结果说明给定起点和短预算能稳定产生可行改进，但均未恢复正式5693.231489\,s水平，不能宣称任意初始解下都稳定得到同等近优解。由于初始即提供可行解，20/20也不是从零构造可行解的成功率。
\fig[.88]{q2_seeds.pdf}{相同起点与150轮预算下20次独立搜索的工期。虚线标出正式方案供参照，但其历史计算预算不同，不用该差值作公平算法排名。}{fig:q2_seed}

\subsubsection{充电时间变化对调度的影响}
将完全充电时间统一乘以$\beta\in\{0.8,0.9,1.0,1.1,1.2,1.3\}$，保留65\%/35\%的两阶段比例。分别检查原始时刻不动的固定日程，以及保持箱组、路线、机型、资源身份和顺序不变、只允许后继开始时刻向后推移的修复日程。对某任务，修复规则为
\begin{equation}
 s_r'=\max\{s_r,t_{\mathrm{aircraft}}',t_{\mathrm{battery}}'\},
\label{eq:q2_repair}
\end{equation}
并按新恢复时间继续传播，不提前任务也不换箱。

充电加快时原日程仍可行，修复规则因不允许提前而不缩短工期。充电时间增加10\%、20\%、30\%时，固定日程分别出现6、6、7处电池接续冲突；顺延修复后均零迟到，工期分别为5802.756723、5955.574471、6108.392219\,s。图\ref{fig:q2_charging}将固定失败与修复代价分开，不能用修复后的成功声称原日程不敏感。
\fig[.88]{q2_charging.pdf}{充电时间倍率与执行代价。叉号表示原固定日程发生资源冲突；曲线为仅顺延开始时刻后的可行结果，不代表重新优化的最短工期。}{fig:q2_charging}

\subsubsection{作业延迟对按时交付的影响}
对23个架次逐一在其第一站交接中增加0、15、30、60、120\,s，每个情景只延迟一个任务，其他交接不变。共115个任务—延迟情景，各检查固定和顺延两种权限，合计230条记录。在当前能耗模型中交接增加时间不另设悬停功率，因而本试验只改变交付、返航和资源时刻，不假称模拟了完整真实悬停能耗。

每个正延迟水平下，固定日程只有9/23个情景保持全部资源与配送条件；其他14个情景首先破坏资源接续，而非出现硬迟到。顺延后23/23均零迟到，所测120\,s单任务延迟的最大工期为5813.231489\,s，比基准增加120\,s。由此可见该方案具有一定时刻修复空间，但原固定日程本身对局部作业延迟敏感。23个确定任务的通过比例不是实际救援成功概率，也不支持多个任务同时延迟的结论。
\fig[.88]{q2_delay.pdf}{单任务首次交接延迟的情景通过数。固定日程与仅顺延修复分开统计；每个延迟水平含23个不同任务情景。}{fig:q2_delay}

\subsubsection{能耗偏差与安全余量的容忍范围}
采用与问题一相同的统一能耗正偏差压力模型，原方案能量边界为
\begin{equation}
 \varepsilon_{\max}=\frac{1-0.2}{1-0.2009738280072748}-1\approx0.121877\%.
\label{eq:q2_energy_limit}
\end{equation}
最低SOC只比下限高0.097383个百分点。本轮在0—2\%范围内对10个能耗水平检查固定与顺延日程，并在边界附近加密。能耗仅增加0.05\%时，原日程已出现3处电池接续冲突；顺延可恢复且不增加最终工期。增加0.13\%后已有一个任务低于20\%SOC，即使调整开始时刻也不能消除能量不可行。这表明“能量容忍范围”与“原时间表容忍范围”仍需分开。

在1\%、2\%压力下进一步开放任务拆分、重新组批、机型与顺序，分别进行250轮搜索。取得的可行备选见表\ref{tab:q2_energy_alternatives}，均由独立扰动物理计算及资源重放核验。它们证明在所测情景下可以通过结构调整恢复，并不证明表中代价最小，也不是原方案已具备同等保障。
\begin{table}[htbp]\centering\small
\caption{能耗偏差下重新调整结构所得可行备选}\label{tab:q2_energy_alternatives}
\begin{tabular}{@{}lrrrcc@{}}
\toprule
偏差 & 架次 & 工期/s & 能耗/kWh & 最低SOC/\% & 独立核验\\
\midrule
1\% & 23 & 5972.072812 & 68.457733 & 20.705 & 通过\\
2\% & 25 & 6650.249809 & 76.308697 & 21.594 & 通过\\
\bottomrule
\end{tabular}
\par\smallskip\parbox{.96\linewidth}{\footnotesize 各250轮搜索；此表为所得可行方案，不宣称情景最优或原方案不变时可用。}
\end{table}


综合而言，名义近优工期与抗扰动能力不是同一个目标。当前日程在小幅时间扰动下可通过顺延恢复，但能量偏差需要更谨慎的结构留量。优先保障稳健执行时，应使用按扰动条件重新验证的备选，而不能仅凭名义最低SOC达标作保证。

\subsubsection{固定资源顺序下的最早顺延修复}
保持箱组、机型、路线、资源身份和顺序不变时，充电或交接时间变化可以沿有向依赖图传播。记任务$r$的原开始为$s_r^0$，同飞机前项为$a(r)$、同电池前项为$b(r)$，扰动后任务与恢复时长为$p'_r,c'_r$，则禁止提前的最早修复满足
\begin{equation}
 s'_r=\max\{s_r^0,\ s'_{a(r)}+p'_{a(r)},\ s'_{b(r)}+p'_{b(r)}+c'_{b(r)}\}.
 \label{eq:q2_maxplus_repair}
\end{equation}
不存在某类前项时省去对应项。原资源顺序来自可行日程，依赖图无环，按拓扑顺序就可完成递推。

该构造在既定顺序和只准向后移动的条件下逐分量最早。初始任务至少要等到原开始和资源可用，取最大值达到下界；若前项已采用最早时刻，任何合法后项都不能早于其释放，所以最大值也给出后项最早开始。沿图归纳即可得到结论。随后再检查箱级截止：如果最早修复已经超期，继续向后移动不能恢复及时性，应改变资源顺序或任务结构。

这使工期变化对应最长依赖路径的增长，而不是扰动总和。两条独立链各推迟30 s可能只延长30 s；同一链上的多个变化则可能累积。单任务延迟实验固定了扰动位置，因此每个恢复结果都可以追溯到具体资源前后关系。

\subsubsection{从压力曲线形成分级调整规则}
充电时间增加30\%时，顺延工期增加415.160730 s，约7.29\%。并行作业吸收了一部分能源恢复延长，只有进入关键接续的部分才传到最后返航。充电加快情景中曲线保持不变，则是因为本次修复不允许提前任务；这与重新优化条件下加快充电可能带来的收益不同。

若仅资源接续被打破而任务仍满足能量安全，优先做时刻修复；若任务自身能耗超过安全预算，改变开始时刻没有作用，应进入重新组批或机型调整。上述两个层次分别对应时间可行性与物理可行性，使压力测试结果可以转为清楚的操作规则。

20次同起点、同预算搜索提供的是解分布，而保存正式方案提供当前最好可行日程，两者在用途上互补。Dolan与Moré提出用明确的性能度量和统一比较对象评估优化软件\cite{dolan2002}。本文相应保留起点、轮数、耗时和全部样本，将短预算稳定性与历史最好解的质量认证分开。研究记录由此既呈现有效改进，也能准确说明计算预算所支持的结论。

